% Gemeinsame Originalaussagen; nur Band 40 registriert sie.
\newcommand{\ReconstructionStatement}[1]{\csname ReconstructionStatement@#1\endcsname}

\expandafter\newcommand\csname ReconstructionStatement@PowerGroupUnitsAreSingletons\endcsname{%
\FormulaThmDeltaK[Einheiten einer Gruppenpotenzhalbgruppe sind genau Einermengen]{%
  \begin{aligned}[t]
    &\GroupAlg{G}{\star}{e}\dsep X\in\PowNE G\vdash{}\\[-2pt]
    &\MonoidUnit{\PowNE G,\star_{\mathcal P},\{e\}}{X}
      \leftrightarrow X\in\operatorname{Sing}(G)
  \end{aligned}
}{PowerGroupUnitsAreSingletons}{%
  \DeltaRow{Mengen}{G\dsep e\dsep X\dsep g}
  \DeltaRow{Funktionen}{\star\dsep\star_{\mathcal P}}
}
}

\expandafter\newcommand\csname ReconstructionStatement@PowerGroupsSingletonLayer\endcsname{%
\FormulaThmDeltaK[Potenzisomorphismen zwischen Gruppen erhalten die Einermengenschicht]{%
  \begin{aligned}[t]
    &\GroupAlg{G}{\star}{e}\dsep\GroupAlg{H}{\diamond}{u}\dsep{}\\[-2pt]
    &\SemigroupIso{F}{\PowNE G,\star_{\mathcal P}}{\PowNE H,\diamond_{\mathcal P}}
      \vdash{}\\[-2pt]
    &F[\operatorname{Sing}(G)]=\operatorname{Sing}(H)
  \end{aligned}
}{PowerGroupsSingletonLayer}{%
  \DeltaRow{Mengen}{G\dsep H\dsep e\dsep u\dsep X}
  \DeltaRow{Funktionen}{F\dsep R=F\restriction_{\operatorname{Sing}(G)}\dsep\star\dsep\diamond}
}
}

\expandafter\newcommand\csname ReconstructionStatement@PowerGroupsSingletonTransport\endcsname{%
\FormulaThmDeltaK[Gruppenrekonstruktion mit vorgegebenen Einermengenbildern]{%
  \begin{aligned}[t]
    &\GroupAlg{G}{\star}{e}\dsep\GroupAlg{H}{\diamond}{u}\dsep{}\\[-2pt]
    &\SemigroupIso{F}{\PowNE G,\star_{\mathcal P}}{\PowNE H,\diamond_{\mathcal P}}
      \vdash{}\\[-2pt]
    &\exists\varphi\,\left(\begin{aligned}[t]
      &\GroupIso{\varphi}{G,\star,e}{H,\diamond,u}\\[-2pt]
      &\land\forall g\in G\,(F(\{g\})=\{\varphi(g)\})
    \end{aligned}\right)
  \end{aligned}
}{PowerGroupsSingletonTransport}{%
  \DeltaRow{Mengen}{G\dsep H\dsep e\dsep u\dsep g}
  \DeltaRow{Funktionen}{F\dsep\varphi\dsep\star\dsep\diamond}
}
}

\expandafter\newcommand\csname ReconstructionStatement@GroupPowerFullCarrierAbsorption\endcsname{%
\FormulaThmDeltaK[Der volle Gruppenträger absorbiert nichtleere Teilmengen]{%
  \begin{aligned}[t]&\GroupAlg{G}{\star}{e}\dsep Y\in\PowNE G\vdash{}\\&Y\star_{\mathcal P}G=G\land G\star_{\mathcal P}Y=G\end{aligned}
}{GroupPowerFullCarrierAbsorption}{%
  \DeltaRow{Mengen}{G\dsep e\dsep X\dsep Y\dsep x\dsep y\dsep z\dsep a\dsep b}
  \DeltaRow{Funktionen}{\star\dsep\star_{\mathcal P}}
}
}

\expandafter\newcommand\csname ReconstructionStatement@SemigroupIsoUnitStabilityTransport\endcsname{%
\FormulaThmDeltaK[Isomorphismen transportieren Einheitenstabilität]{%
  \begin{aligned}[t]&\MonoidAlg{M}{\star}{e}\dsep\MonoidAlg{N}{\diamond}{d}\dsep{}\\&\SemigroupIso{h}{M,\star}{N,\diamond}\dsep\operatorname{UStab}_{M,\star,e}(z)\vdash{}\\&\operatorname{UStab}_{N,\diamond,d}(h(z))\end{aligned}
}{SemigroupIsoUnitStabilityTransport}{%
  \DeltaRow{Mengen}{M\dsep N\dsep e\dsep d\dsep z\dsep u\dsep w}
  \DeltaRow{Funktionen}{\star\dsep\diamond\dsep h}
}
}

\expandafter\newcommand\csname ReconstructionStatement@PowerMonoidUnitSetStable\endcsname{%
\FormulaThmDeltaK[Die Einheitenmenge ist im Potenzmonoid einheitenstabil]{%
  \MonoidAlg{A}{\star}{e}\vdash \operatorname{UStab}_{\PowNE A,\star_{\mathcal P},\{e\}}(A^\times)
}{PowerMonoidUnitSetStable}{%
  \DeltaRow{Mengen}{A\dsep e\dsep u\dsep W}
  \DeltaRow{Funktionen}{\star\dsep\star_{\mathcal P}}
}
}

\expandafter\newcommand\csname ReconstructionStatement@PowerMonoidUnitStableAbsorption\endcsname{%
\FormulaThmDeltaK[Einheitenstabilität liefert Absorption durch die Einheitenmenge]{%
  \begin{aligned}[t]&\MonoidAlg{A}{\star}{e}\dsep X\in\PowNE A\dsep{}\\&\operatorname{UStab}_{\PowNE A,\star_{\mathcal P},\{e\}}(X)\vdash{}\\&A^\times\star_{\mathcal P}X=X\land X\star_{\mathcal P}A^\times=X\end{aligned}
}{PowerMonoidUnitStableAbsorption}{%
  \DeltaRow{Mengen}{A\dsep e\dsep X\dsep x\dsep z\dsep u}
  \DeltaRow{Funktionen}{\star\dsep\star_{\mathcal P}}

  \DeltaRow{Abkürzungen}{U=A^\times\dsep P=\PowNE A}
}
}

\expandafter\newcommand\csname ReconstructionStatement@PowerMonoidUnitSubsetCriterion\endcsname{%
\FormulaThmDeltaK[Produktkriterium für Teilmengen der Einheitenmenge]{%
  \begin{aligned}[t]&\MonoidAlg{A}{\star}{e}\dsep Y\in\PowNE A\vdash{}\\&Y\in\PowNE{A^\times}\leftrightarrow Y\star_{\mathcal P}A^\times=A^\times\end{aligned}
}{PowerMonoidUnitSubsetCriterion}{%
  \DeltaRow{Mengen}{A\dsep e\dsep Y\dsep y}
  \DeltaRow{Funktionen}{\star\dsep\star_{\mathcal P}}
}
}

\expandafter\newcommand\csname ReconstructionStatement@LargePowerUnitSetImage\endcsname{%
\FormulaThmDeltaK[Bild der gesamten Einheitenmenge]{%
  \begin{aligned}[t]&\MonoidAlg{A}{\star}{e_A}\dsep\MonoidAlg{B}{\diamond}{e_B}\dsep{}\\&\SemigroupIso{F}{\PowNE A,\star_{\mathcal P}}{\PowNE B,\diamond_{\mathcal P}}\vdash F(A^\times)=B^\times\end{aligned}
}{LargePowerUnitSetImage}{%
  \DeltaRow{Mengen}{A\dsep B\dsep e_A\dsep e_B}
  \DeltaRow{Funktionen}{\star\dsep\diamond\dsep\star_{\mathcal P}\dsep\diamond_{\mathcal P}\dsep F}

  \DeltaRow{Abkürzungen}{U=A^\times\dsep V=B^\times}
  \DeltaRow{}{P=\PowNE A\dsep Q=\PowNE B}
}
}

\expandafter\newcommand\csname ReconstructionStatement@LargePowerUnitSubsetImage\endcsname{%
\FormulaThmDeltaK[Bilder nichtleerer Teilmengen der Einheitenmenge]{%
  \begin{aligned}[t]&\MonoidAlg{A}{\star}{e_A}\dsep\MonoidAlg{B}{\diamond}{e_B}\dsep\SemigroupIso{F}{\PowNE A,\star_{\mathcal P}}{\PowNE B,\diamond_{\mathcal P}}\dsep{}\\&F(A^\times)=B^\times\dsep Y\in\PowNE{A^\times}\vdash F(Y)\in\PowNE{B^\times}\end{aligned}
}{LargePowerUnitSubsetImage}{%
  \DeltaRow{Mengen}{A\dsep B\dsep e_A\dsep e_B\dsep Y}
  \DeltaRow{Funktionen}{\star\dsep\diamond\dsep\star_{\mathcal P}\dsep\diamond_{\mathcal P}\dsep F}
}
}

\expandafter\newcommand\csname ReconstructionStatement@LargePowerSemigroupUnitGroupReconstruction\endcsname{%
\FormulaThmDeltaK[Große Potenzhalbgruppen rekonstruieren die
Einheitengruppen]{%
  \begin{aligned}[t]
    &\MonoidAlg{A}{\star}{e_A}\dsep
      \MonoidAlg{B}{\diamond}{e_B}\dsep{}\\[-2pt]
    &\SemigroupIso{F}
      {\PowNE A,\star_{\mathcal P}}
      {\PowNE B,\diamond_{\mathcal P}}\vdash{}\\[-2pt]
    &F(A^\times)=B^\times
      \land F[\PowNE{A^\times}]=\PowNE{B^\times}
  \end{aligned}%
}{LargePowerSemigroupUnitGroupReconstruction}{%
  \DeltaRow{Mengen}{A\dsep B\dsep A^\times\dsep B^\times
    \dsep X\dsep Y\dsep e_A\dsep e_B}
  \DeltaRow{Funktionen}{\star\dsep\diamond\dsep\star_{\mathcal P}
    \dsep\diamond_{\mathcal P}\dsep F}

  \DeltaRow{Abkürzungen}{U=A^\times\dsep V=B^\times}
  \DeltaRow{}{P=\PowNE A\dsep Q=\PowNE B}
  \DeltaRow{}{R=F\restriction_{\PowNE U}}
}
}

\expandafter\newcommand\csname ReconstructionStatement@GroupPowerSemigroupRigidity\endcsname{%
\FormulaThmDeltaK[Einseitige Gruppenstarrheit der großen
Potenzhalbgruppe]{%
  \begin{aligned}[t]
    &\GroupAlg{A}{\star}{e_A}\dsep
      \SemigroupAlg{B}{\diamond}\dsep{}\\[-2pt]
    &\SemigroupIso{F}
      {\PowNE A,\star_{\mathcal P}}
      {\PowNE B,\diamond_{\mathcal P}}\vdash{}\\[-2pt]
    &\exists e_B\in B\,\exists f\,
      \bigl(
        \GroupAlg{B}{\diamond}{e_B}\land
        \GroupIso{f}{A,\star,e_A}{B,\diamond,e_B}
      \bigr)
  \end{aligned}%
}{GroupPowerSemigroupRigidity}{%
  \DeltaRow{Mengen}{A\dsep B\dsep A^\times\dsep B^\times
    \dsep E\dsep X\dsep Y\dsep e_A\dsep e_B}
  \DeltaRow{Funktionen}{\star\dsep\diamond\dsep\star_{\mathcal P}
    \dsep\diamond_{\mathcal P}\dsep F\dsep f}

  \DeltaRow{Abkürzungen}{U=A^\times\dsep V=B^\times}
  \DeltaRow{}{P=\PowNE A\dsep Q=\PowNE B}
}
}

% Der intern benannte Produktzeuge behält seine ursprüngliche Hilfssatznummer.
% Die formale Ableitung steht unverändert in der ausgelagerten Absorptionstabelle.
\newcommand{\ReconstructionProductMemberFormula}{%
  \begin{aligned}[t]
    &\SemigroupAlg{G}{\star}\dsep X,Y\in\PowNE G\dsep x\in X\dsep y\in Y\\[-2pt]
    &\vdash x\star y\in X\star_{\mathcal P}Y
  \end{aligned}%
}
\newcommand{\ReconstructionProductMemberStatement}{%
  \par\Needspace{8\baselineskip}%
  \begingroup
  \setcounter{proofpartnr}{0}%
  \ProofPartSetParent{\theformulaThm}%
  \refstepcounter{proofpartnr}%
  \edef\mxPPLbl{thm:pp:\mxProofPartParent:\arabic{proofpartnr}}%
  \label{\mxPPLbl}%
  \ThmLookupRegisterId{GroupPowerFullCarrierAbsorptionProductMember}{theorem}{\mxPPLbl}%
  \noindent\textbf{Produktzeuge}\hfill\ThmLookupEmitRef{theorem}{\mxPPLbl}\par
  \[\ReconstructionProductMemberFormula\]
  \noindent Die Ableitung bildet den ersten Teil der ausgelagerten Tabelle
  (\ReconstructionProofLink{GroupPowerFullCarrierAbsorption}).\par
  \endgroup
}
